% ============================================================================
% Section 3: Related Work
% ============================================================================

\section{Related Work}\label{sec:related_work}

We survey several lines of research that inform the design of our benchmark, improvements made to our expert harness (\pokeagent{}), and AutoEvolve algorithm. These range from existing contributions to experiential learning, recursive language models, and various self-adaptive methods.

\subsection{Experiential Learning}\label{sec:experiential_learning}

A growing body of work explores how agents can accumulate and leverage experience over extended interactions, moving beyond static policy deployment toward systems that learn through interaction. While experiential learning has a deep tradition in the RL paradigm through self-play and policy optimization across episodes, experiential learning in the LLM setting emphasizes structured knowledge accumulation within and across sessions, using the model's own reasoning and reflection as the primary learning signal.

AriGraph~\citep{anokhin2024arigraph} constructs knowledge graph world models with episodic memory for LLM agents, enabling them to maintain structured representations of their environment across interactions. This approach addresses the inability of standard LLM agents to retain and organize discoveries, spatial relationships, and causal knowledge beyond the current context window. AriGraph demonstrates that persistent, structured memory can compensate for the finite context window and enable agents to more effectively reason over accumulated experience.

Language Models Meet World Models~\citep{xiang2023embodied} showed that grounding LLMs in embodied experience, rather than relying solely on pretraining knowledge, significantly improves their performance on downstream reasoning tasks. Their method finetunes on structured experiential data collected during embodied play, enabling agents to adapt to novel environments that diverge from their training distribution.

Recent work on in-context reinforcement learning (ICRL) has demonstrated that LLMs can perform implicit policy improvement from in-context reward signals~\citep{song2025icrl, karten2025economist}. We reason that if LMs can implicitly improve policies from reward signals presented in-context, then explicitly grounding these signals in the agent's trajectory and reflection process may further complement this capability. This principle informs several of the methods discussed in Chapter~\ref{sec:autoevolve}.

\subsection{Skill Learning}\label{sec:skill_learning}

A recurring theme in experiential learning is the curation of reusable behavioral skills. Voyager~\citep{wang2023voyager} demonstrated that language model agents can progressively acquire capabilities through interaction, storing successful programs in a growing code library indexed by natural language descriptions.

SkillRL~\citep{xia2026skillrl} provides a comprehensive treatment of this process. Rather than accumulating raw, often noisy and redundant trajectories, SkillRL introduces an experience-based distillation mechanism that extracts high-level behavioral patterns into a hierarchical skill library (SkillBank). An adaptive retrieval strategy enables the agent to access both general and task-specific heuristics at inference time, while a recursive evolution mechanism allows the skill library to co-evolve with the agent's policy during reinforcement learning training. On the ALFWorld, WebShop, and seven search-augmented benchmarks, SkillRL achieves state-of-the-art performance, outperforming strong baselines by over 15.3\% while significantly reducing the token footprint. The connections between SkillRL's distillation approach and our harness's skill-learning mechanism are revisited in Chapters~\ref{sec:pokeagent} and~\ref{sec:autoevolve}.

\subsection{Agent Memory}\label{sec:agent_memory}

Memory is the primary mechanism by which experiential learning persists beyond the immediate context window, and its design has emerged as a critical architectural decision for long-horizon agents. Agent memory systems generally organize knowledge along three axes: \emph{episodic memory} (logs of specific interactions and events), \emph{semantic memory} (accumulated facts, strategies, and world knowledge), and \emph{procedural memory} (learned skills and behavioral routines).

Generative Agents~\citep{park2023generative} introduced an influential memory architecture that combines recency, importance, and relevance-based retrieval to select which memories to surface during reasoning. This design enables believable long-term behavior in simulated social environments and has influenced subsequent work on agent memory, including retrieval-augmented approaches that treat memory as an external knowledge base to be queried at inference time.

MemGPT~\citep{packer2023memgpt} draws an explicit analogy between LLM context management and operating-system virtual memory. By implementing a hierarchical memory system with fast (in-context) and slow (external storage) tiers, MemGPT enables agents to intelligently move information in and out of their active context window, maintaining coherent behavior across interactions that would otherwise exceed context limits. This approach is particularly relevant to the long-context settings we consider, where the space of all possible information that might be relevant to an agent exceeds what can be reasonably stored at each step. Our expert harness leverages this hierarchical memory abstraction to efficiently store and dynamically access memory across all three axes.

\subsection{Recursive Language Models}\label{sec:rlm}

Recursive Language Models (RLMs)~\citep{zhang2025rlm} programmatically interact with and recurse over their own context. Rather than treating the context window as a passive repository of information, RLMs treat it as an active workspace that the model can read from, write to, and restructure. RLMs are particularly relevant to the problem of \emph{context rot}, where LMs experience a degradation in coherent reasoning as context windows grow large. As agents accumulate interaction history, the ratio of relevant to irrelevant information in-context deteriorates, leading to repetitive behavior, lost goals, and contradictory reasoning. These failure modes are documented across long-horizon agent systems from Claude Plays Pok\'{e}mon~\citep{anthropic2025visible} to Gemini Plays Pok\'{e}mon~\citep{zhang2025geminiplayspokemon}. The RLM approach to context-management directly informs the design of both our expert harness (Chapter~\ref{sec:pokeagent}), which maintains a structured trajectory window that subagents process and summarize, and our AutoEvolve algorithm (Chapter~\ref{sec:autoevolve}), which can be understood as a form of recursive self-modification operating over trajectory segments.

\subsection{Self-Adaptive Methods}\label{sec:self_adaptation}

Several lines of work explore how language models and LLM-based agents can improve their own performance through self-directed learning, making different tradeoffs between scope of adaptation and computational overhead.

\paragraph{Prompt Evolution.} GEPA~\citep{agrawal2025gepa} uses genetic evolution over prompt candidates, guided by natural-language reflections on agent trajectories. It runs an agent multiple times, collects reasoning traces and outcomes, and has the model analyze these traces to propose prompt edits. In line with the principles of natural selection, candidates that improve performance are retained for the next generation. GEPA demonstrated that this in-context, reflective evolution can outperform RL-based fine-tuning across several benchmarks, establishing prompt evolution as a viable alternative to weight updates. The method's primary limitation, for our purposes, is its requirement of full episodic completion and environment resets between iterations, making it less directly applicable to dynamic long-horizon settings where resets are expensive or unavailable. AutoEvolve (Chapter~\ref{sec:autoevolve}) addresses this limitation directly by using the global session log of the agents reasoning traces and environment interactions as an implicit learning signal without environment reset.

\paragraph{LLM-as-a-Judge.} The use of language models as automated evaluators has become widespread, with MT-Bench and Chatbot Arena~\citep{zheng2023mtbench} establishing the paradigm of LLM-as-a-Judge for open-ended evaluation. Subsequent work has identified systematic biases in this approach, including position bias and verbosity bias~\citep{shi2024judgingthejudges}, and has extended the paradigm to multimodal settings~\citep{chen2024mllmjudge}. AutoEvolve's evolutionary optimization relies on an LLM judge module that evaluates trajectory segments and produces structured critiques. We address known biases through careful prompt design and by grounding the judge's evaluations in concrete game-state metrics (milestone progress, step counts, spatial positioning) in addition to qualitative assessment.
